\chapter{Algorithmic Design of the Simulator}
\label{c:algorithmic}

%%%%%%%%%%%%%%%%%%%%%%%%%%%%%%%%%%%%%%%%%%%%%%%%%%%%%%%%%%%%%%%%%%%%%%%%%%%%%%%%

% ==========================================
% Section: Matrix-Free Direct State Evolution
% ==========================================
\section{Matrix-Free Direct State Evolution}\label{sec:matrix-free-direct-state-evolution}
The core computational challenge in simulating many-body quantum systems lies in the efficient application of excitation operators to the underlying state vector. Constructing and applying global unitary matrices is highly memory-prohibitive for large Hilbert spaces. Therefore, the simulator leverages a direct state-vector approach, exploiting the algebraic properties of the excitation generator to execute updates directly on the computational basis.

To translate this theoretical framework into a high-performance algorithmic routine, we must deconstruct the evaluation of the unitary operation defined in Equation \eqref{unitary_operation}. As an initial step, we isolate the fundamental mechanics by considering the qubit excitation framework. This isolated model provides a computational baseline that is then subsequently generalized to incorporate fermionic anticommutation relations.

\subsection{Qubit Baseline}\label{subsec:qubit-baseline}
Rather than relying on computationally prohibitive full matrix exponentiation to evaluate the time-evolution operator $U = \exp(\sum \theta \hat{E})$, the direct state evolution approach exploits the algebraic properties of the generator to execute localized state updates. To compute the excitation efficiently, we must simplify the evaluation of the unitary operation defined in Equation \eqref{unitary_operation}. As an initial simplification, we consider the qubit excitation framework, which can later be generalized to account for fermionic differences. 

Consider an initial quantum state upon which the excitation is applied. Recalling the generalized generator from Equation \eqref{eq:mult_qu_exc_generator}, the operation is fully defined by the orbital sets $\boldsymbol{k}$ and $\boldsymbol{l}$, expressing which electrons are exchanged. With these target orbitals and the initial state, the algorithm possesses all the necessary parameters. Given the generator $G^{\boldsymbol{k}}_{\boldsymbol{l}}$, the unitary operator can be expanded as:
\begin{equation}
	U(\theta) = e^{-i\frac{\theta}{2}G^{\boldsymbol{k}}_{\boldsymbol{l}}} = \cos\left(\frac{\theta}{2}\right) I - i\sin\left(\frac{\theta}{2}\right)G^{\boldsymbol{k}}_{\boldsymbol{l}} + \left(1 - \cos\left(\frac{\theta}{2}\right)\right) P_0.
\end{equation}
Here, $P_0$ represents the projector onto the nullspace of the generator $G^{\boldsymbol{k}}_{\boldsymbol{l}}$. This nullspace is spanned by all computational basis states $\ket{\boldsymbol{i}}$ for which $G^{\boldsymbol{k}}_{\boldsymbol{l}} \ket{\boldsymbol{i}} = 0$.

When applied to a specific electronic configuration $\ket{\Phi}$, the unitary transformation acts trivially (as the identity operator) if the configuration resides within the nullspace of the generator. Otherwise, it induces a rotation between the original configuration and the $n$-fold excited configuration:
\begin{equation} \label{eq:qubit-excitation-unitary-transformation-configuration-differences}
	\qubitExcitationUnitaryTransformationConfigurationDifferences
\end{equation}

To implement the active transformation algorithmically, we examine the two components of the rotational term. The first term, governed by $\cos\left(\frac{\theta}{2}\right) I$, is computationally straightforward. It is expressed as a direct scaling of the initial state's probability amplitude, which can be expressed algorithmically as an update:
\begin{equation}
	c_{\boldsymbol{i}} \leftarrow c_{\boldsymbol{i}} \cos\left(\frac{\theta}{2}\right).\label{eq:qubit-excitation-update-basis-state}
\end{equation}

The second component of the rotation governs the coupling and is more complex as it involves the explicit action of the generator. The term $-i\sin\left(\frac{\theta}{2}\right)G^{\boldsymbol{k}}_{\boldsymbol{l}}$ can be decomposed into two distinct operational steps. First, the algorithm computes the partner state $\ket{\boldsymbol{i'}}$ resulting from the generator's application onto the basis state:
\begin{align}
	c \ket{\boldsymbol{i'}} &= -iG^{\boldsymbol{k}}_{\boldsymbol{l}} \ket{\boldsymbol{i}}, \label{eq:qubit-excitation-baseline-partner-state-1} \\
	-iG^{\boldsymbol{k}}_{\boldsymbol{l}} &= \left( \prod_n \sigma^-_{k_n} \sigma^+_{l_n} - \text{h.c.} \right). \label{eq:qubit-excitation-baseline-partner-state-2}
\end{align}
Here, the partner state $\ket{\boldsymbol{i'}}$ represents the specific electronic configuration generated by the excitation, which is also introduced into the superposition. By absorbing the factor of $-i$ into the generator's evaluation, the transitional coefficient $c$ is restricted to purely real values and even $c = \pm 1$. The sign of $c$ depends strictly on the directionality of the particle exchange (i.e., whether the operation acts as an excitation or de-excitation on the specific state).

Finally, this term induces an amplitude update on the newly acquired partner state. The corresponding probability amplitude is updated as:
\begin{equation}
	c_{\boldsymbol{i'}} \leftarrow c_{\boldsymbol{i'}} + c \sin\left(\frac{\theta}{2}\right). \label{eq:qubit-excitation-partner-state-update-phase}
\end{equation}
\cite[Chapter II.A]{kottmann2020feasible}

% ==========================================
% Subsection: Incorporating the Fermionic Phase
% ==========================================
\subsection{Incorporating the Fermionic Phase}\label{subsec:incorporating-the-fermionic-phase}

While the foundational baseline algorithm effectively handles qubit-based rotations, it is inherently insufficient for simulating the characteristics of true fermionic many-body systems. To accurately model fermions, the computational routine must replace the qubit excitation generator with its fermionic counterpart, the multi-excitation generator $G_{\boldsymbol{pq}}$, as defined in Equation \eqref{eq:mult_fe_exc_generator}.

Within a classical simulation framework, direct access to the state vector enables the evaluation of fermionic excitations without mapping to qubit Pauli operators. By representing computational basis states as Fock space occupation vectors $\ket{\mathbf{n}} = \ket{n_0, n_1, \dots, n_{M-1}}$ with $n_i \in \{0, 1\}$, the algebraic action of second-quantized fermionic creation ($a_k^\dagger$) and annihilation ($a_k$) operators can be defined directly on the basis states.

The fermionic anticommutation relations dictate that acting on a target spin-orbital $k$ introduces a parity phase determined by the total occupation of all preceding orbitals in the canonical ordering. Algebraically, the action of these operators on a Fock state $\ket{\mathbf{n}}$ is expressed as:
\begin{align} \label{eq:annhiliation-creation-algebraic-definition}
	a_k \ket{\mathbf{n}} &= \delta_{n_k, 1} (-1)^{\sum_{j=0}^{k-1} n_j} \ket{n_0, \dots, n_k - 1, \dots, n_{M-1}}, \\
	a_k^\dagger \ket{\mathbf{n}} &= \delta_{n_k, 0} (-1)^{\sum_{j=0}^{k-1} n_j} \ket{n_0, \dots, n_k + 1, \dots, n_{M-1}},
\end{align}
where $\delta$ denotes the Kronecker delta, enforcing the Pauli exclusion principle. The phase factor $(-1)^{\sum_{j=0}^{k-1} n_j}$ evaluates to $+1$ for an even preceding occupation count and $-1$ for an odd count.

To evaluate an arbitrary $n$-fold excitation, the constituent operators are applied sequentially. For each operator targeting index $k$, the parity phase is computed against the instantaneous state vector. By updating the occupation value $n_k \mapsto 1 - n_k$ immediately after each valid operator application, this algebraic formulation tracks all necessary anticommutation phases natively, eliminating the need for Jordan-Wigner parity strings.

This fundamental difference in dynamic behavior enforces the strict antisymmetry required of fermionic wavefunctions, originating from the Pauli exclusion principle. Algorithmically, it can be recognized that the nullspace projector $P_0$ acts purely based on the occupancy of the initial and target orbitals, ignoring intermediate phase dynamics. Therefore, the correct evaluation of the active state rotation relies entirely on accurately tracking and applying this classically computed parity phase.

Finally, note that the distinct computational advantage of this classical approach. Implementing the fermionic phase switch on physical quantum hardware is notoriously complex; applying long chains of Pauli-Z operators requires deep CNOT ladders, which introduce significant gate overhead and noise. By determining the phase algebraically through direct state occupation counting, our simulator bypasses this physical bottleneck.

\cite[chapter III.B]{anand2022quantum}
\cite[chapter 2.1.1]{schleich2025quantum}

% ==========================================
% Section: Fermionic SWAP (fSWAP)
% ==========================================
\section{Fermionic SWAP (fSWAP)}
Routing and reordering fermionic orbitals requires the fermionic SWAP (fSWAP) gate to preserve anticommutation relations. In the computational basis, this gate is defined as:
\begin{equation}
	U_{\text{fSWAP}} = \begin{pmatrix}
						   1 & 0 & 0 & 0 \\
						   0 & 0 & 1 & 0 \\
						   0 & 1 & 0 & 0 \\
						   0 & 0 & 0 & -1
	\end{pmatrix}
\end{equation}
The negative phase on the $|11\rangle$ element enforces the required exchange sign when swapping two occupied orbitals.

\cite[chapter II.]{Verstraete2009quantum}

% ==========================================
% Section: Expectation Value Evaluation
% ==========================================
\section{Expectation Value Evaluation}\label{sec:expectation-value-evaluation}
Extracting physical observables requires evaluating transition matrix elements $\bra{\phi} H \ket{\psi}$. Decomposing the Hamiltonian into a linear combination of fermionic operator strings $H = \sum_k w_k \hat{O}_k$, and expanding the state vectors in the $2^N$-dimensional computational basis ($\ket{\psi} = \sum_x c_x \ket{x}$ and $\bra{\phi} = \sum_y d_y^* \bra{y}$), yields for each operator term $\hat{O}_k$:
\begin{equation}
	\bra{\phi} \hat{O}_k \ket{\psi} = \sum_{x} \sum_{y} d_y^* c_x \bra{y} \hat{O}_k \ket{x}
\end{equation}
Applying a fermionic operator string $\hat{O}_k$ to a basis state $\ket{x}$ either annihilates the state, if the operation violates Pauli exclusion, or maps it to a single new basis state $\ket{x'_k}$ with an accumulated parity phase $p_{k,x} \in \{-1, 1\}$, such that $\hat{O}_k \ket{x} = p_{k,x} \ket{x'_k}$. Due to basis orthonormality, the inner product evaluates to a Kronecker delta:
\begin{equation}
	\bra{y} \ket{x'_k} = \delta_{y, x'_k} = \begin{cases} 1 & \text{if } y = x'_k \\ 0 & \text{if } y \neq x'_k \end{cases}
\end{equation}
This orthogonality collapses the double summation over basis states into a single linear sum over the non-zero amplitudes of $\ket{\psi}$. The total expectation value is thus computed as:
\begin{equation}
	\bra{\phi} H \ket{\psi} = \sum_k w_k \sum_{x} d_{x'_k}^* p_{k,x} c_x\label{eq:expectation-value-decomposed}
\end{equation}
Algorithmatically, resolving $\ket{x'_k}$ via bitwise manipulations allows the simulator to evaluate observables without constructing dense matrix operators.
\cite[chapter II.A, II.B]{kohda2022quantum}